\section*{Nutzung künstlicher Intelligenzen}
In der folgenden Tabelle ist aufgeführt, welche KI-Werkzeuge in dieser Arbeit wofür genutzt wurden. 
Alle Ergebnisse der KI-Werkzeuge wurden von der Gruppe geprüft, im Labor getestet und bei Bedarf angepasst.

\begin{tabular}{|L{2.6cm}|L{9.1cm}|L{3.5cm}|}
  \hline
  \textbf{Werkzeug} & \textbf{Beschreibung der Nutzung} & \textbf{Arbeitsschritt / Kapitel} \\
  \hline
  \multicolumn{3}{|l|}{\textbf{1. Allgemein}} \\
  \hline
  ChatGPT, Claude & Verschaffen eines Überblicks über die Aufgabenstellung und Planung des Vorgehens. & Planung \\
  \hline
  Claude & Unterstützung bei der Ausarbeitung des Bedien- und Anzeigekonzepts (Belegung der Taster, Bedeutung der LEDs). & Konzept, Kapitel~\ref{sec:programmkonzept} \\
  \hline
  \multicolumn{3}{|l|}{\textbf{2. Programmierung im TIA Portal}} \\
  \hline
  ChatGPT, Claude & Unterstützung beim Entwurf des Zustandsautomaten und beim Erstellen des SCL-Codes für den Baustein \texttt{TempDisplay\_Ctrl}. & Programmierung \\
  \hline
  ChatGPT, Claude & Unterstützung bei der Fehlersuche während der Inbetriebnahme, z.\,B. bei Referenzfahrt und Drehrichtung. & Test im Labor \\
  \hline
  \multicolumn{3}{|l|}{\textbf{3. Erstellung der Dokumentation}} \\
  \hline
  Claude & Unterstützung bei Gliederung, Formulierung und sprachlicher Überarbeitung der Kapitel sowie beim Erstellen der Ablaufdiagramme in LaTeX. & Dokumentation \\
  \hline
  GitHub Copilot & Unterstützung beim Verfassen einzelner Formulierungen und bei der Korrektur von Rechtschreibfehlern. & Dokumentation \\
  \hline
  ChatGPT & Unterstützung bei LaTeX und bei der Analyse von Compilerfehlern. & Dokumentation \\
  \hline
\end{tabular}
